\begin{ouverture}[college1]{1}{Matière : mesurer, séparer, identifier}{Partie 1 · Organisation et transformations de la matière}{Comment distinguer deux liquides transparents ?}
\objectif{Décrire états, mélanges et changements d’état.}
\objectif{Exploiter une masse volumique et des mesures.}
\objectif{Choisir une méthode de séparation et contrôler sa conclusion.}
\prerequis{Masse en g ou kg ; volume en mL, L ou m³.}
\prerequis{Proportionnalité et division.}
\end{ouverture}
\pcPage{Le cours}{Matière : mesurer, séparer, identifier}
\phantomsection\label{ess-1}
\essSection{1}{États et modèle microscopique}
Un solide a une forme et un volume propres ; un liquide a un volume propre et prend la forme du récipient ; un gaz occupe tout le volume disponible. Dans le modèle, les entités sont proches dans un solide ou un liquide, très espacées dans un gaz. Une fusion, une solidification, une vaporisation ou une liquéfaction change leur organisation, sans former une nouvelle espèce. La sublimation et la condensation solide relient directement solide et gaz.\par
\pcFigure{figures/b01-etats.pdf}{Même espèce ; organisation différente. Les particules ne sont pas à l’échelle.}
\essSection{2}{Mesurer un changement d’état}
À pression fixée, un corps pur présente une température de changement d’état déterminée. Sur une courbe, un palier peut correspondre à la coexistence de deux états, alors que de l’énergie est transférée. Dans un système fermé sans perte de matière, la masse se conserve ; le volume peut changer. Un palier seul ne prouve pas toute la composition d’un échantillon : il faut connaître le protocole et sa précision.\par
\clearpage
\essSection{3}{Masse volumique}
La masse volumique est $\rho=m/V$, soit $m=\rho V$ et $V=m/\rho$. En unités SI : m en kg, V en m³, $\rho$ en kg/m³. On peut aussi employer g et mL, ou g et cm³, si les unités sont cohérentes. $1\unite{mL}=1\unite{cm^3}$ et $1\unite{g/mL}=1\,000\unite{kg/m^3}$. La température, les bulles et une tare incorrecte influencent la mesure.\par
\pcFigure{figures/b01-rho.pdf}{Pour un matériau homogène à température fixée, masse et volume sont proportionnels.}
\essSection{4}{Corps purs, mélanges et séparation}
Un corps pur contient une seule espèce chimique ; un mélange en contient plusieurs. Homogène signifie qu’on ne distingue pas plusieurs phases à l’œil nu. Deux liquides miscibles forment une phase. Une filtration retient des particules solides, mais pas le sel dissous ; une distillation permet de recueillir un liquide vaporisé puis liquéfié. La solubilité est une limite de dissolution à température donnée. Une eau limpide n’est pas nécessairement potable.\par
\begin{pcEncadre}{method}{Un exemple expliqué}Un solide de $156\unite g$ déplace l’eau de 40 à $60\unite{mL}$. Son volume est $20\unite{cm^3}$ ; $\rho=156/20=7{,}8\unite{g/cm^3}$. C’est compatible avec le fer, sans identification unique. Il faut une immersion complète, sans dissolution ni bulle.\end{pcEncadre}
\pcPage{Expérimenter et retenir}{Matière : mesurer, séparer, identifier}
\phantomsection\label{activite-1}
\begin{pcEncadre}{method}{Investigation en classe}En classe, mesurer la masse de trois volumes d’un même liquide avec tare. Tracer m en fonction de V au tableur, comparer les quotients et discuter les écarts. Concevoir ensuite un protocole distinct pour un solide insoluble. Indiquer ce qui est mesuré directement, calculé et supposé.\end{pcEncadre}
\pcSousTitre{Ce que je dois savoir expliquer}\begin{questions}
\item Distinguer espèce, état et mélange.
\item Expliquer conservation de masse et variation possible du volume.
\item Choisir des unités cohérentes pour $
ho=m/V$.
\end{questions}
\begin{pcEncadre}{attention}{Une erreur à éviter}La masse volumique n’est pas la masse : doubler le volume d’un matériau homogène double sa masse, sans doubler sa masse volumique.\end{pcEncadre}
\pcRepere{cor-activite-1}{Pistes pour l’investigation}\pcGrille
\clearpage\markright{Je démarre}
\rubrique{Je démarre}
\par\noindent\begin{minipage}{\linewidth}
\begin{exercice}[1]{Trois états}
\begin{questions}
\item Nommer l’état qui conserve sa forme.
\item Quel état remplit tout un récipient fermé ?
\end{questions}
\end{exercice}
\end{minipage}
\par\noindent\begin{minipage}{\linewidth}
\begin{exercice}[1]{Une eau salée}
\begin{questions}
\item Corps pur ou mélange ?
\item Homogène signifie-t-il pur ?
\end{questions}
\end{exercice}
\end{minipage}
\par\noindent\begin{minipage}{\linewidth}
\begin{exercice}[1]{Les bons mots}
\begin{questions}
\item Nommer le passage liquide-gaz.
\item Nommer le passage gaz-liquide.
\end{questions}
\end{exercice}
\end{minipage}
\par\noindent\begin{minipage}{\linewidth}
\begin{exercice}[1]{Volumes}
\begin{questions}
\item Convertir $0{,}35\unite L$ en mL.
\item Convertir $18\unite{cm^3}$ en mL.
\end{questions}
\end{exercice}
\end{minipage}
\par\noindent\begin{minipage}{\linewidth}
\begin{exercice}[1]{Le quotient}
\begin{questions}
\item Calculer $\rho$ pour 40 g et 50 mL.
\item Donner le résultat en kg/m³.
\end{questions}
\end{exercice}
\end{minipage}
\clearpage\markright{Je m’entraîne}
\rubrique{Je m’entraîne}
\par\noindent\begin{minipage}{\linewidth}
\begin{exercice}[2]{Solide immergé}
\pcFigureEx{figures/rev-ex-immersion.pdf}{}
\begin{questions}
\item Une éprouvette passe de 28 à 43 mL. Trouver le volume du solide.
\item Sa masse est 40,5 g. Calculer sa masse volumique.
\end{questions}
\end{exercice}
\end{minipage}
\par\noindent\begin{minipage}{\linewidth}
\begin{exercice}[2]{Une masse prévue}
\begin{questions}
\item Une huile a $\rho=0{,}90\unite{g/mL}$. Quelle masse pour 250 mL ?
\item La même huile chauffée occupe plus de volume à masse conservée. Comment évolue $\rho$ ?
\end{questions}
\end{exercice}
\end{minipage}
\par\noindent\begin{minipage}{\linewidth}
\begin{exercice}[2]{Retrouver un volume}
\begin{questions}
\item Un morceau de cuivre de 178 g a $\rho=8{,}9\unite{g/cm^3}$. Trouver V.
\item Pourquoi une pièce creuse ne suit-elle pas le même calcul avec son volume extérieur ?
\end{questions}
\end{exercice}
\end{minipage}
\par\noindent\begin{minipage}{\linewidth}
\begin{exercice}[2]{Une limite de dissolution}
\begin{questions}
\item À température fixée, 100 mL d’eau dissolvent au plus 36 g de sel. Que devient un ajout de 40 g ?
\item Estimer la masse maximale pour 250 mL d’eau dans les mêmes conditions.
\end{questions}
\end{exercice}
\end{minipage}
\par\noindent\begin{minipage}{\linewidth}
\begin{exercice}[2]{Interpréter un palier}
\pcFigureEx{figures/rev-ex-ebullition.pdf}{Courbe illustrant le palier de trois minutes décrit dans la question.}
\begin{questions}
\item Pendant trois minutes, une eau pure bout vers 100 °C à pression usuelle. Pourquoi la température ne monte-t-elle pas malgré le chauffage ?
\item Sa masse baisse dans le bécher ouvert. La conservation est-elle violée ?
\end{questions}
\end{exercice}
\end{minipage}
\clearpage\markright{Je raisonne}
\rubrique{Je raisonne}
\par\noindent\begin{minipage}{\linewidth}
\begin{exercice}[3]{Choisir une séparation}
\begin{questions}
\item Un mélange contient eau, sel dissous et sable. Proposer deux opérations pour recueillir de l’eau séparée.
\item La première opération suffit-elle à prouver la potabilité ?
\end{questions}
\end{exercice}
\end{minipage}
\par\noindent\begin{minipage}{\linewidth}
\begin{exercice}[3]{Une mesure discutable}
\begin{questions}
\item Un bloc de 54 g paraît déplacer 25 mL, alors que son volume réel est 20 mL. Quel $\rho$ calcule-t-on ?
\item Une bulle accrochée peut-elle expliquer le sens de l’erreur ?
\end{questions}
\end{exercice}
\end{minipage}
\par\noindent\begin{minipage}{\linewidth}
\begin{exercice}[3]{Identifier avec une plage}
\begin{questions}
\item Une mesure donne $2{,}68\unite{g/cm^3}$ avec une incertitude de lecture estimée à $0{,}05\unite{g/cm^3}$. Comparer à 2,70 pour l’aluminium.
\item Peut-on conclure « aluminium pur certain » ?
\end{questions}
\end{exercice}
\end{minipage}
\par\noindent\begin{minipage}{\linewidth}
\begin{exercice}[3]{Glace et océans}
\begin{questions}
\item Expliquer pourquoi la fonte d’une glace continentale peut augmenter le niveau marin.
\item Citer un autre mécanisme physique lié au réchauffement.
\end{questions}
\end{exercice}
\end{minipage}
\par\noindent\begin{minipage}{\linewidth}
\begin{exercice}[3]{Construire une preuve}
\begin{questions}
\item Deux liquides ont le même aspect. Proposer une comparaison quantitative sans les goûter.
\item Pourquoi conserver température et protocole comparables ?
\end{questions}
\end{exercice}
\end{minipage}
\pcPage{Défi du chapitre}{Matière : mesurer, séparer, identifier}
\phantomsection\label{defi-ex-1}
\pcDefiTitre{1}{Le flacon inconnu}
\begin{questions}
\item Un flacon vide pèse 45 g ; avec 80 mL de liquide, 109 g. Déterminer $\rho$.
\item Références : eau 1,00 ; huile 0,90 ; éthanol 0,79 g/mL. Quelle hypothèse est la plus compatible ?
\item Proposer une limite et une amélioration du protocole.
\end{questions}
\pcGrille
